\chapter{Complex Analysis}
\label{ch:complex}

\medskip

\textbf{Notation.} $U\subset\mathbb{C}$ open; $\Delta_r(a)=\{z:|z-a|<r\}$; $\partial U$ --- boundary.
$H(U)$ --- holomorphic (complex differentiable) functions on $U$; $f^{(n)}$ --- $n$-th derivative.
$\oint_\gamma f\,dz$ --- contour integral; $\operatorname{ind}_\gamma(a)=\frac1{2\pi i}\oint_\gamma\frac{dz}{z-a}$ --- winding number.
$\operatorname{Res}(f,a)$ --- residue at $a$.
A curve $\gamma$ is \textit{null-homotopic} in $U$ if it can be continuously contracted to a point within $U$.
An \textit{entire} function is holomorphic on all of $\mathbb{C}$.
An \textit{annulus} is a region $r<|z-a|<R$.

\medskip

\section{Cauchy--Riemann equations}
\label{sec:cauchy-riemann}

$f=u+iv$ is holomorphic (complex differentiable) on $U$ $\iff$ $u,v$ are $C^1$ and satisfy
$$u_x=v_y,\quad u_y=-v_x.$$
Equivalently $\partial_{\overline z}f=0$. A holomorphic function is infinitely differentiable and analytic.

\medskip

\section{Cauchy's integral theorem and formula}
\label{sec:cauchy-integral}

\subsection{Cauchy's theorem}
\label{subsec:cauchy-theorem}

$f\in H(U)$, $\gamma$ is a null-homotopic closed curve in $U$ $\Rightarrow$ $\oint_\gamma f(z)\,dz=0$.

\subsection{Cauchy's integral formula}
\label{subsec:cauchy-formula}

$f\in H(U)$, $\overline{\Delta_r(a)}\subset U$. For $z\in\Delta_r(a)$:
$$f(z)=\frac1{2\pi i}\oint_{|\zeta-a|=r}\frac{f(\zeta)}{\zeta-z}\,d\zeta.$$
Consequence (derivatives): $f^{(n)}(a)=\frac{n!}{2\pi i}\oint\frac{f(\zeta)}{(\zeta-a)^{n+1}}\,d\zeta$.

\subsection{Cauchy estimate and Liouville}
\label{subsec:cauchy-estimate}

$|f^{(n)}(a)|\le\frac{n!}{r^n}\max_{|\zeta-a|=r}|f(\zeta)|$.
\textit{Liouville:} a bounded entire function is constant. \textit{Fundamental theorem of algebra:} every non-constant polynomial has a root in $\mathbb{C}$.

\medskip

\section{Power series and analyticity}
\label{sec:power-series-complex}

\subsection{Taylor series}
\label{subsec:taylor-series-complex}

Holomorphic $f$ is analytic: $f(z)=\sum_{n\ge0}a_n(z-a)^n$, converging in $\Delta_R(a)$ where $R=$ distance to the nearest singularity. $a_n=f^{(n)}(a)/n!$.

\subsection{Laurent series}
\label{subsec:laurent-series}

$f$ holomorphic on an annulus $r<|z-a|<R$:
$$f(z)=\sum_{n=-\infty}^\infty c_n(z-a)^n,\quad c_n=\frac1{2\pi i}\oint_{|\zeta-a|=\rho}\frac{f(\zeta)}{(\zeta-a)^{n+1}}\,d\zeta.$$

\medskip

\section{Zeros and singularities}
\label{sec:singularities}

\subsection{Zeros}
\label{subsec:zeros-complex}

$f\not\equiv0$ holomorphic, $f(a)=0$ $\Rightarrow$ $f(z)=(z-a)^m g(z)$, $g(a)\neq0$. Zeros are isolated unless $f\equiv0$.

\textit{Identity theorem:} $f=g$ on a set with a limit point in $U$ $\Rightarrow$ $f\equiv g$ on $U$.

\subsection{Classification of isolated singularities}
\label{subsec:singularities-class}

\noindent
\begin{itemize}
\item \textit{Removable:} Laurent series has no negative terms; $\lim_{z\to a}f(z)$ finite.
\item \textit{Pole:} finite negative part; $\lim_{z\to a}f(z)=\infty$.
\item \textit{Essential:} infinitely many negative terms; $\lim_{z\to a}f(z)$ does not exist.
\end{itemize}

\subsubsection{Pole criterion}
\label{subsubsec:pole-criterion}

If $|f(z)|\to\infty$ as $z\to a$, then $a$ is a pole. $a$ is a pole $\iff$ $1/f$ has a removable zero.

\medskip

\section{Residue theorem}
\label{sec:residue-theorem}

\subsection{Residue}
\label{subsec:residue}

$\operatorname{Res}(f,a)=c_{-1}$ for a Laurent expansion at $a$.
\begin{itemize}
\item Simple pole: $\operatorname{Res}(f,a)=\lim_{z\to a}(z-a)f(z)=\dfrac{g(a)}{h'(a)}$ if $f=g/h$, $h(a)=0$, $h'(a)\neq0$.
\item Pole of order $m$: $\operatorname{Res}(f,a)=\frac1{(m-1)!}\lim_{z\to a}\frac{d^{m-1}}{dz^{m-1}}\bigl((z-a)^m f(z)\bigr)$.
\end{itemize}

\subsection{Residue theorem}
\label{subsec:residue-theorem}

$f\in H(U\setminus S)$, $\gamma$ null-homotopic in $U$, avoids $S$, finite $S\cap\operatorname{int}\gamma$:
$$\oint_\gamma f(z)\,dz=2\pi i\sum_{a\in S}\operatorname{Res}(f,a)\operatorname{ind}_\gamma(a).$$

\textit{When to use.} Computing real integrals $\int_{-\infty}^\infty R(x)\,dx$, $\int_0^{2\pi} R(\cos\theta,\sin\theta)\,d\theta$, sums $\sum_{n\in\mathbb{Z}} f(n)$ via $\pi\cot\pi z$.

\medskip

\section{Maximum modulus and Schwarz lemma}
\label{sec:maximum-modulus}

\subsection{Maximum modulus principle}
\label{subsec:maximum-modulus}

$f\in H(U)$, $|f|$ attains a local maximum at $a\in U$ $\Rightarrow$ $f$ is constant on the component. Minimum principle: if $f\neq0$, $|f|$ attains a minimum $\Rightarrow$ $f$ constant.

\textit{Corollary (Hadamard's three circles):} $\ln M(r)$ is a convex function of $\ln r$, where $M(r)=\max_{|z|=r}|f(z)|$.

\subsection{Schwarz lemma}
\label{subsec:schwarz-lemma}

$f:\Delta_1(0)\to\Delta_1(0)$, $f(0)=0$ $\Rightarrow$ $|f(z)|\le|z|$ and $|f'(0)|\le1$. Equality at $z\neq0$ or $|f'(0)|=1$ $\Rightarrow$ $f(z)=e^{i\theta}z$.
